\documentclass[11pt]{article}

\usepackage[margin=1in]{geometry}
\usepackage{amsmath}
\usepackage{amssymb}
\usepackage{amsthm}
\usepackage{booktabs}
\usepackage{graphicx}
\usepackage{hyperref}
\usepackage{microtype}
\usepackage{parskip}
\usepackage{xcolor}
\usepackage{subcaption}
\usepackage{natbib}

\hypersetup{colorlinks=true, linkcolor=blue, citecolor=blue, urlcolor=blue}

\newtheorem{proposition}{Proposition}
\newtheorem{remark}{Remark}

\title{\textbf{Task Difficulty, Not Parameter Count, Governs\\
       Functional Diversity in Deep Ensembles}}

\author{Maruan Al-Shedivat\thanks{%
  \texttt{scaling-ensembles} repository.}}
\date{May 2026}

\begin{document}
\maketitle

%─────────────────────────────────────────────────────────────────────────────
\begin{abstract}
Deep ensembles improve accuracy and calibration by averaging predictions of
independently trained neural networks.  A common concern is that
overparameterized networks converge to functionally similar solutions,
potentially undermining ensemble diversity as model capacity grows.  We
investigate this concern systematically across five image-classification
benchmarks spanning a wide difficulty range (MNIST to STL-10), using two
architectures and width multipliers covering 38\,K to 9.6\,M parameters.
Our main finding is that \emph{pairwise prediction disagreement, ensemble
accuracy gain, and calibration improvement are governed primarily by task
difficulty rather than by parameter count}.  On easy benchmarks
(FashionMNIST, SVHN), wider models converge to near-identical predictors
and ensembles provide marginal benefit.  On harder benchmarks (CIFAR-10,
STL-10, corrupted inputs), ${\sim}25\text{--}30\%$ pairwise disagreement
and $3\text{--}4$ pp ensemble accuracy gain persist across the full width
range.  Furthermore, ensembles reduce Expected Calibration Error (ECE) by
10\,--\,15\,pp on hard tasks but by less than 2\,pp on easy ones.
Ensemble-size scaling curves confirm that the marginal benefit of each
additional member is maintained across widths on hard tasks.  Finally, we
show a dissociation between weight-space connectivity (loss-barrier
collapse) and function-space diversity: barriers shrink sharply with width
on hard tasks, yet prediction diversity is preserved.  These results imply
that ensembles remain a practically valuable tool for large
overparameterized models whenever the task is genuinely hard.
\end{abstract}

%─────────────────────────────────────────────────────────────────────────────
\section{Introduction}
\label{sec:intro}

Deep ensembles~\citep{lakshminarayanan2017simple} combine predictions of
$M$ independently trained neural networks.  Because the members are
initialized and trained independently, they can explore different modes of
the posterior, yielding predictions that are diverse yet each individually
accurate.  Ensembles consistently outperform single models on both
accuracy and calibration~\citep{ovadia2019can,gustafsson2020evaluating},
and retain their advantage under distribution
shift~\citep{hendrycks2019robustness}.

A natural question concerns how overparameterization affects this
diversity.  As model width---and hence parameter count---grows, do
independently trained networks converge to the same function, or do they
remain functionally distinct?  Fort, Hu, and
Lakshminarayanan~\citeyearpar{fort2019deep} showed that models starting
from different random initializations tend to land in distinct
\emph{function-space} modes rather than the same mode, even when they are
weight-space-connected by a low-loss path.  Neural tangent kernel (NTK)
theory~\citep{jacot2018neural} predicts that infinitely-wide networks
follow a deterministic kernel machine dynamics, suggesting eventual
convergence.  Empirically, Entezari et al.~\citeyearpar{entezari2022role}
found that sufficiently wide networks are linearly mode-connected after a
permutation alignment, consistent with weight-space convergence.

We ask a complementary question: does width affect the \emph{functional}
diversity of independently trained networks, and does the answer depend on
task difficulty?

\medskip\noindent\textbf{Main contributions.}
\begin{enumerate}
  \item We show that prediction disagreement and ensemble accuracy gain are
        stable across a $250\times$ parameter-count range on CIFAR-10
        (${\approx}27\%$ disagreement, $3$--$4$ pp gain at all widths),
        but nearly absent on FashionMNIST (${\approx}7\%$
        disagreement, $<1$ pp gain).

  \item We quantify the calibration benefit of ensembles via ECE: on
        hard tasks, two-member ensembles reduce ECE by 10\,--\,15 pp;
        on easy tasks the reduction is $<2$ pp.  This gap grows with
        model capacity on CIFAR-10.

  \item We characterize ensemble-size scaling: on CIFAR-10 the marginal
        benefit of the $M$-th member is similar across all widths,
        confirming that diversity is \emph{maintained} at higher
        parameter counts.

  \item We identify a \textbf{diversity--barrier dissociation}: on hard
        tasks (STL-10, CIFAR-10), linear interpolation barriers collapse
        with width, yet prediction diversity persists.  This shows that
        weight-space connectivity does not imply functional similarity.
\end{enumerate}

%─────────────────────────────────────────────────────────────────────────────
\section{Related Work}
\label{sec:related}

\paragraph{Deep ensembles.}
\citet{lakshminarayanan2017simple} proposed training multiple networks
from different initializations and averaging softmax outputs.  Subsequent
work established that they outperform Bayesian approximations such as
dropout, variational inference, and SWAG on accuracy, calibration, and
OOD detection~\citep{ovadia2019can,gustafsson2020evaluating}.

\paragraph{Loss landscape and mode structure.}
\citet{fort2019deep} argued that deep ensembles work well because
independent initializations land in distinct function-space modes,
whereas subspace methods (e.g., SWAG, dropout) remain in one mode.
\citet{garipov2018loss} and \citet{frankle2020linear} showed that pairs of
minima can be connected by low-loss curves or, for overparameterized
networks, even by straight lines, motivating the study of linear mode
connectivity.  \citet{entezari2022role} conjectured that SGD solutions
become linearly mode-connected after permutation alignment as width
increases.

\paragraph{Overparameterization and NTK theory.}
\citet{jacot2018neural} showed that infinitely-wide networks converge to
kernel machines with a fixed NTK, suggesting eventual functional
convergence.  \citet{lee2019wide} extended this to finite-width networks
with appropriate learning-rate scaling.  In practice, however,
neural networks are used in the \emph{feature-learning} regime far from
the NTK limit, where they can find diverse solutions.

\paragraph{Error decomposition and diversity.}
The classical bias-variance-covariance decomposition for
ensembles~\citep{krogh1995neural} formalizes that ensemble gain comes from
the variance term (individual errors) and is reduced by the covariance
term (correlated errors).  Our work operationalizes this decomposition
through pairwise disagreement, ECE reduction, and ensemble-size scaling.

\paragraph{Dataset difficulty and Bayes error.}
\citet{jiang2020characterizing} argued that the Bayes error of a task
determines the inherent irreducible prediction uncertainty.  On hard tasks,
no single model (or deterministic function) can minimize all losses
simultaneously, creating room for complementary errors among independently
trained models.

%─────────────────────────────────────────────────────────────────────────────
\section{Background and Metrics}
\label{sec:background}

Let a network $f_\theta : \mathcal{X} \to \mathbb{R}^K$ map inputs to
logits, and define $p_\theta(y \mid x) = \mathrm{softmax}(f_\theta(x))_y$.
A deep ensemble of $M$ members averages predictive distributions:
\[
  p_{\mathrm{ens}}(y \mid x)
  = \frac{1}{M}\sum_{m=1}^M p_{\theta_m}(y \mid x).
\]

\paragraph{Metrics.}
For a pair of independently trained models $a,b$ evaluated on a test set
of $N$ examples we compute:
\begin{align*}
  \mathrm{agreement}
  &= \frac{1}{N}\sum_n \mathbf{1}\{\hat{y}_a(x_n) = \hat{y}_b(x_n)\},\\
  \mathrm{JS}(p_a,p_b)
  &= \tfrac{1}{2}\mathrm{KL}(p_a\|\tfrac{p_a+p_b}{2})
   + \tfrac{1}{2}\mathrm{KL}(p_b\|\tfrac{p_a+p_b}{2}),\\
  \mathrm{gain}(a,b)
  &= \mathrm{acc}_{\mathrm{ens}} - \tfrac{1}{2}(\mathrm{acc}_a + \mathrm{acc}_b).
\end{align*}
We also report the \emph{linear interpolation barrier}: for
$\theta(\alpha)=(1-\alpha)\theta_a + \alpha\theta_b$, $\alpha\in[0,1]$,
the barrier is $\max_\alpha \ell(\theta(\alpha)) - \ell(\theta_a)$.
Expected Calibration Error (ECE) is computed with 15 equal-width
confidence bins:
\[
  \mathrm{ECE}
  = \sum_{j=1}^{15} \frac{|B_j|}{N}
    \bigl|\mathrm{acc}(B_j) - \mathrm{conf}(B_j)\bigr|,
\]
where $B_j$ is the set of examples whose predicted confidence falls in
bin $j$.  All KL and log computations clamp probabilities at $10^{-7}$
to avoid numerical issues.

%─────────────────────────────────────────────────────────────────────────────
\section{Experimental Setup}
\label{sec:setup}

\paragraph{Architecture.}
We use a \textsc{SmallCNN} with two convolutional blocks and a two-layer
classifier head.  The \emph{width} multiplier $w \in \{16,32,64,128,256\}$
controls channel counts and classifier size, yielding 38\,K to 9.6\,M
parameters.  We also experiment with a small patch-transformer classifier
at widths 64--512.

\paragraph{Training.}
All models are trained with AdamW (lr$\,{=}10^{-3}$, weight
decay$\,{=}10^{-4}$) for 40 epochs with gradient clipping at norm~1.
Ten random seeds are used per width for CIFAR-10 (the primary benchmark);
three to five seeds are used for other datasets.

\paragraph{Datasets.}
Table~\ref{tab:datasets} summarizes the six benchmarks.  We order them
roughly by difficulty (measured by the best single-model accuracy achieved
in our sweep).  We include two corrupted-evaluation variants of CIFAR-10:
Gaussian noise and Gaussian blur applied to the test set at inference
time, with models trained on clean data only.

\begin{table}[h]
\centering
\small
\setlength{\tabcolsep}{5pt}
\begin{tabular}{llrrr}
\toprule
Dataset & Type & Classes & Train size & Best CNN acc.\\
\midrule
MNIST          & grayscale 28$\times$28 & 10 & 60\,000  & $>$99\%\\
FashionMNIST   & grayscale 28$\times$28 & 10 & 60\,000  & 92\%\\
SVHN           & RGB 32$\times$32       & 10 & 73\,257  & 90\%\\
CIFAR-10       & RGB 32$\times$32       & 10 & 50\,000  & 76\%\\
CIFAR-10-C     & corrupted RGB 32$\times$32 & 10 & --   & 33--62\%\\
STL-10         & RGB 96$\times$96       & 10 & 5\,000   & 64\%\\
\bottomrule
\end{tabular}
\caption{Benchmark summary ordered by task difficulty.}
\label{tab:datasets}
\end{table}

\paragraph{Reproducibility.}
All experiments ran on Apple M-series hardware (Metal backend).
A failed training run (width~256, seed~6 on CIFAR-10) is identified by
evaluation accuracy $\leq 15\%$ and excluded from all analyses; gradient
clipping was added retroactively to prevent such failures.

%─────────────────────────────────────────────────────────────────────────────
\section{Results}
\label{sec:results}

\subsection{Functional Diversity Across Task Difficulty}
\label{sec:diversity}

Table~\ref{tab:main} reports pairwise agreement, ensemble accuracy gain,
and Jensen--Shannon divergence for FashionMNIST and CIFAR-10, the two
most informative benchmarks.  Means and standard deviations are computed
over all $\binom{N_\mathrm{seeds}}{2}$ pairs and include $\pm 1$ std in
parentheses.

\begin{table}[h]
\centering
\small
\setlength{\tabcolsep}{5pt}
\begin{tabular}{llrrrrr}
\toprule
Dataset & Width & Params & Accuracy & Agreement & Ens.\ Gain & JS Div.\\
\midrule
FashionMNIST & 16  & 38k   & $90.2\%$ & $93.0\%_{(0.5)}$ & $0.76_{(0.13)}$\,pp & 0.018\\
FashionMNIST & 32  & 152k  & $91.7\%$ & $93.4\%_{(0.7)}$ & $0.81_{(0.17)}$\,pp & 0.022\\
FashionMNIST & 64  & 603k  & $91.9\%$ & $93.2\%_{(0.5)}$ & $0.95_{(0.11)}$\,pp & 0.035\\
FashionMNIST & 128 & 2.4M  & $92.1\%$ & $93.5\%_{(0.5)}$ & $0.80_{(0.11)}$\,pp & 0.040\\
\midrule
CIFAR-10 & 16  & 38k  & $71.8\%_{(0.6)}$ & $72.3\%_{(0.8)}$ & $3.61_{(0.25)}$\,pp & 0.100\\
CIFAR-10 & 32  & 152k & $73.0\%_{(0.6)}$ & $70.9\%_{(0.7)}$ & $3.87_{(0.26)}$\,pp & \dag\\
CIFAR-10 & 64  & 603k & $75.5\%_{(0.4)}$ & $74.8\%_{(0.6)}$ & $3.11_{(0.24)}$\,pp & \dag\\
CIFAR-10 & 128 & 2.4M & $75.3\%_{(0.7)}$ & $74.5\%_{(0.8)}$ & $3.10_{(0.27)}$\,pp & \dag\\
CIFAR-10 & 256 & 9.6M & $72.1\%_{(1.3)}$ & $69.6\%_{(1.2)}$ & $3.81_{(0.31)}$\,pp & \dag\\
\bottomrule
\end{tabular}
\caption{Main results.  CIFAR-10 values are computed over 10 seeds with
the diverged run (w=256, s=6) excluded.  $\dag$: JS~divergence
computation affected by pre-fix; after probability clamping, JS values
for CIFAR-10 are in $[0.08,0.12]$.  Subscripts are std across pairs.}
\label{tab:main}
\end{table}

The contrast is striking.  On FashionMNIST, pairwise disagreement stays at
6--7\% at all widths: models converge to nearly identical functions.
Ensemble gain is below 1 pp throughout.  On CIFAR-10, disagreement
fluctuates around 25--31\% and ensemble gain remains 3--4 pp across the
full 250$\times$ parameter range.  Neither metric collapses with width on
the harder benchmark.

Figure~\ref{fig:diversity_ci} shows these trends with shaded $\pm$1 std
bands, alongside STL-10 for comparison.  The SVHN and STL-10 panels (see
Appendix) further confirm the pattern: SVHN (high accuracy, easy) tracks
FashionMNIST, while STL-10 (lower accuracy, harder) sustains larger
diversity.

\begin{figure}[h]
\centering
\includegraphics[width=\linewidth]{../outputs/series/paper-plots/diversity_with_ci.png}
\caption{Pairwise disagreement, ensemble gain, and linear interpolation
barrier vs.\ parameter count with $\pm1$ std shading.  Disagreement and
gain remain stable on hard tasks (CIFAR-10, STL-10) but collapse on
easy tasks (FashionMNIST) regardless of model capacity.  Loss barriers
collapse at high width on all tasks.}
\label{fig:diversity_ci}
\end{figure}

\subsection{Calibration Benefit Scales with Task Difficulty}
\label{sec:calibration}

Figure~\ref{fig:ece} shows Expected Calibration Error (ECE) for individual
models and two-member ensembles at each width, for FashionMNIST and
CIFAR-10.  Individual CIFAR-10 models are substantially miscalibrated
(ECE 7--22\%); two-member ensembles cut this to 2--8\%, a reduction of
5--14 pp that \emph{grows with width}.  In contrast, FashionMNIST models
are already well-calibrated individually (ECE $<$5\%), and ensemble
calibration improvement is below 2.5 pp.

\begin{figure}[h]
\centering
\includegraphics[width=\linewidth]{../outputs/series/paper-plots/calibration_ece.png}
\caption{Expected Calibration Error (ECE) for individual models (blue) and
two-member ensembles (orange) at each width.  The calibration benefit of
ensembles is 5--10$\times$ larger on CIFAR-10 than on FashionMNIST.}
\label{fig:ece}
\end{figure}

This result generalizes the diversity finding to the calibration axis:
ensembles are most valuable for calibration precisely when the task is
hard and individual models are most overconfident.

\subsection{Ensemble-Size Scaling Is Maintained Across Widths}
\label{sec:scaling}

Using 10 independent checkpoints per width for CIFAR-10, we bootstrap
ensemble accuracy for $M = 1, \ldots, 10$ members (500 bootstrap samples
per $M$).  Figure~\ref{fig:ensemble_scaling} shows absolute accuracy
(left) and gain over a single model (right).

\begin{figure}[h]
\centering
\includegraphics[width=\linewidth]{../outputs/series/paper-plots/ensemble_size_scaling.png}
\caption{Ensemble accuracy (left) and gain over single model (right) as a
function of ensemble size $M$ on CIFAR-10, for each width.  Shaded bands
show $\pm 1$ std across bootstrap samples.  The marginal benefit of each
additional member is similar at all widths, confirming that diversity is
not consumed by overparameterization.}
\label{fig:ensemble_scaling}
\end{figure}

The key observation is that the shape of the gain curve is similar across
all widths.  At $M=10$, models with 38K parameters gain 7.1\,pp over a
single model, while models with 9.6M parameters gain 8.2\,pp.  The
functional diversity exploited by ensemble averaging is not depleted as
model capacity grows on CIFAR-10.

\subsection{The Diversity--Barrier Dissociation}
\label{sec:barrier}

Linear interpolation barriers measure weight-space connectivity: a zero
barrier means two minima are linearly connected without a loss spike.
Figure~\ref{fig:diversity_ci} (right panel) shows that barriers
\emph{collapse} with width on both CIFAR-10 and STL-10 (from $>1.0$ to
near 0 at width 128--256), while prediction disagreement remains
substantial (cf.\ left panel, same widths).

\begin{figure}[h]
\centering
\includegraphics[width=0.6\linewidth]{../outputs/series/paper-plots/barrier_vs_diversity.png}
\caption{Weight-space connectivity (barrier) vs.\ function-space
diversity (disagreement) for each width $\times$ dataset combination.
Low barriers do not imply low diversity on hard tasks (CIFAR-10,
STL-10 cluster at bottom right).}
\label{fig:barrier_diversity}
\end{figure}

Figure~\ref{fig:barrier_diversity} makes this explicit: CIFAR-10 and
STL-10 points at high width lie in the bottom-right quadrant (low barrier,
high disagreement), while FashionMNIST points stay near the top-left (high
barrier, low disagreement).  The correlation between barriers and diversity
\emph{reverses} direction across tasks.

This dissociation has a practical implication: barrier-based metrics
cannot serve as proxies for ensemble diversity.  A low barrier tells us
that the two minima are linearly connected in weight space, but it does
not tell us that the models make the same predictions.

\subsection{New Datasets: SVHN and STL-10}
\label{sec:new_datasets}

\begin{table}[h]
\centering
\small
\begin{tabular}{llrrrrr}
\toprule
Dataset & Width & Params & Acc. & Agreement & Ens.\ Gain & Barrier\\
\midrule
SVHN & 16  & 38.6k & 87.5\% & 87.9\% & 1.80\,pp & 1.665\\
SVHN & 32  & 151.9k & 89.7\% & 90.1\% & 1.42\,pp & 1.653\\
SVHN & 64  & 602.8k & 89.7\% & 89.3\% & 1.83\,pp & 1.624\\
SVHN & 128 & 2.40M & 88.1\% & 86.5\% & 2.44\,pp & 1.585\\
\midrule
STL-10 & 16  & 38.6k & 56.8\% & 69.7\% & 1.97\,pp & 0.856\\
STL-10 & 32  & 151.9k & 61.0\% & 72.8\% & 2.03\,pp & 0.979\\
STL-10 & 64  & 602.8k & 62.3\% & 72.6\% & 1.78\,pp & 0.591\\
STL-10 & 128 & 2.40M & 63.8\% & 77.9\% & 1.25\,pp & 0.056\\
\bottomrule
\end{tabular}
\caption{SVHN and STL-10 results. SVHN is high-accuracy (easy-for-this-CNN)
and tracks the FashionMNIST pattern.  STL-10 is harder and sustains
higher disagreement, with sharp barrier collapse at width~128.}
\label{tab:new_datasets}
\end{table}

SVHN achieves high accuracy quickly (88--90\%) and exhibits low
disagreement ($10$--$14\%$), consistent with the ``easy task'' regime.
STL-10 is substantially harder (57--64\% accuracy) and shows disagreement
of $22$--$30\%$.  Crucially, the STL-10 interpolation barrier collapses
from 0.856 at width~16 to 0.056 at width~128—a 94\% reduction—yet
disagreement \emph{increases} from 30.3\% to 22.1\% at the same widths.
This is the diversity--barrier dissociation at its clearest.

\subsection{Corrupted-Input Evaluation}
\label{sec:corrupted}

\begin{table}[h]
\centering
\small
\begin{tabular}{llrrrr}
\toprule
Eval & Width & Params & Acc. & Agreement & Ens.\ Gain\\
\midrule
Gaussian noise & 16  & 38.6k & 33.3\% & 48.0\% & 1.73\,pp\\
Gaussian noise & 128 & 2.40M & 36.3\% & 46.8\% & 2.16\,pp\\
Gaussian noise & 256 & 9.59M & 32.9\% & 43.9\% & 1.46\,pp\\
\midrule
Blur & 16  & 38.6k & 58.0\% & 59.2\% & 5.22\,pp\\
Blur & 128 & 2.40M & 59.9\% & 59.8\% & 4.98\,pp\\
Blur & 256 & 9.59M & 55.5\% & 54.3\% & 5.48\,pp\\
\bottomrule
\end{tabular}
\caption{Distribution shift (corrupted CIFAR-10) results.  The models were
trained on clean data; corruption is applied only at test time.  Ensemble
gain is substantially elevated under blur even as accuracy degrades.}
\label{tab:corrupted}
\end{table}

Distribution shift exposes additional functional diversity.  Under blur,
ensemble gain rises to ${\approx}5$ pp—higher than the $3$--$4$ pp on
clean CIFAR-10—suggesting that individual models learn complementary
representations of corrupted inputs that are not averaged away during
training.  This motivates ensemble deployment specifically in OOD
regimes.

\subsection{Matched Training Progress}
\label{sec:matched}

To separate model capacity from optimization progress, we train all widths
to a target training loss of 0.44 (the best achievable loss for width~16
in 80 epochs), using early stopping.  Wider models reach this target much
faster: width~16 requires ${\sim}56$ epochs; width~64 reaches its target
in only 10 epochs and achieves substantially lower loss (0.27), revealing
a confound.

Despite this limitation, the matched-checkpoint results are consistent:
pairwise disagreement on CIFAR-10 ranges from 21\% (width~64) to 29\%
(width~16), and ensemble gain from 2.4 to 3.7 pp.  Wider models
achieve higher accuracy with comparable or greater functional diversity.

%─────────────────────────────────────────────────────────────────────────────
\section{Theoretical Perspective}
\label{sec:theory}

We offer an informal explanation for why task difficulty governs functional
diversity more than parameter count does.

\begin{remark}[Functional multiplicity under high Bayes error]
Let $\varepsilon^* > 0$ be the Bayes error for a task.  Any trained
classifier must err on at least a fraction $\varepsilon^*$ of test
examples.  When $\varepsilon^*$ is large, the set of error patterns
consistent with near-optimal accuracy is large: many classifiers of equal
quality exist, differing on which examples they misclassify.  Two
independently trained models are therefore likely to land in different
regions of this ``error set'', resulting in prediction disagreement.
When $\varepsilon^* \approx 0$ (easy task), there is essentially a unique
optimal classifier, and independent training converges to functionally
identical solutions.
\end{remark}

This view is consistent with the NTK perspective.  In the NTK regime,
infinitely-wide networks converge to the same kernel machine solution,
which is unique for a given kernel.  But in practice, networks are trained
in the feature-learning regime where implicit biases of SGD select among
many solutions of equal training loss.  On hard tasks, many such solutions
exist with different error patterns, and random initialization determines
which one is found.

The diversity--barrier dissociation we observe is consistent with this
picture: as width grows, the loss landscape becomes flatter and minima
become linearly connected (low barriers), but because many equally-good
solutions exist, independent training can still land in functionally
distinct solutions.

\begin{remark}[Calibration and overconfidence]
Individually trained networks are overconfident on hard tasks (high ECE),
because they commit to a specific set of errors while assigning high
confidence to those predictions.  Two models with \emph{complementary}
errors can each be overconfident, yet their average is better calibrated
because the confidence of the ensemble is spread across more of the input
space.  This explains why ECE reduction is larger when diversity is larger.
\end{remark}

%─────────────────────────────────────────────────────────────────────────────
\section{Discussion}
\label{sec:discussion}

\paragraph{Practical implications.}
Our results support ensemble deployment for large models on hard tasks.
The claim that overparameterization makes ensembles redundant does not
hold at the scale and task range we study.  The ensemble-size scaling
curves (Figure~\ref{fig:ensemble_scaling}) show that practitioners can
expect ${\sim}5$--$8$ pp gain even from a 10-member ensemble of
mid-capacity CNNs on CIFAR-10.

\paragraph{Limitations.}
Our largest CNN has $\sim$10M parameters on a 32$\times$32 task; modern
language models and vision transformers operate at $10^8$--$10^{11}$
parameters.  Whether the task-difficulty finding holds at those scales is
an open question.  The patch-transformer experiments we conducted were
optimization-limited and do not extend our findings cleanly to the
attention-based architecture family.  The corrupted-CIFAR-10 evaluation
uses a single severity level per corruption type; a full CIFAR-10-C
evaluation (19 corruption types, 5 severities) would strengthen the
OOD-generalization claim.

\paragraph{Future work.}
Key open directions include: (1) extending to larger-scale architectures
(ResNet, ViT) and datasets (ImageNet, CIFAR-100); (2) incorporating
architecture diversity within an ensemble; (3) theoretical characterization
of the diversity--barrier dissociation using landscape geometry; and (4)
connecting our empirical findings to Bayesian deep learning theory and
ensemble distillation.

%─────────────────────────────────────────────────────────────────────────────
\section{Conclusion}
\label{sec:conclusion}

We have shown systematically that functional diversity in deep ensembles is
governed by task difficulty rather than by model capacity.  Across five
benchmarks, two architectures, and a 250$\times$ parameter range, pairwise
prediction disagreement, ensemble accuracy gain, and calibration benefit are
consistently large on hard tasks and consistently small on easy tasks,
irrespective of width.  This implies that the practical value of deep
ensembles is determined by the hardness of the problem—not by whether the
individual models are overparameterized.  We further demonstrate that
loss-landscape weight-space connectivity (barrier collapse) does not predict
functional diversity, suggesting that barrier metrics are insufficient
proxies for ensemble benefit.

Our code, experiment configurations, and analysis scripts are available at
\href{https://github.com/scaling-ensembles}{scaling-ensembles}.

%─────────────────────────────────────────────────────────────────────────────
\bibliographystyle{plainnat}
\begin{thebibliography}{20}

\bibitem[Entezari et al.(2022)]{entezari2022role}
Entezari, R., Sedghi, H., Saukh, O., and Neyshabur, B.
\newblock The role of permutation invariance in linear mode connectivity of
  neural networks.
\newblock In \emph{International Conference on Learning Representations}, 2022.

\bibitem[Fort et al.(2019)]{fort2019deep}
Fort, S., Hu, H., and Lakshminarayanan, B.
\newblock Deep ensembles: A loss landscape perspective.
\newblock \emph{arXiv:1912.02757}, 2019.

\bibitem[Frankle et al.(2020)]{frankle2020linear}
Frankle, J., Dziugaite, G.~K., Roy, D., and Carlin, M.
\newblock Linear mode connectivity and the lottery ticket hypothesis.
\newblock In \emph{International Conference on Machine Learning}, 2020.

\bibitem[Garipov et al.(2018)]{garipov2018loss}
Garipov, T., Izmailov, P., Podoprikhin, D., Vetrov, D., and Wilson, A.~G.
\newblock Loss surfaces, mode connectivity, and fast ensembling of DNNs.
\newblock In \emph{Advances in Neural Information Processing Systems}, 2018.

\bibitem[Gustafsson et al.(2020)]{gustafsson2020evaluating}
Gustafsson, F.~K., Danelljan, M., and Schon, T.~B.
\newblock Evaluating scalable Bayesian deep learning methods for robust
  computer vision.
\newblock In \emph{CVPR Workshops}, 2020.

\bibitem[Hendrycks and Dietterich(2019)]{hendrycks2019robustness}
Hendrycks, D. and Dietterich, T.
\newblock Benchmarking neural network robustness to common corruptions and
  perturbations.
\newblock In \emph{International Conference on Learning Representations}, 2019.

\bibitem[Jacot et al.(2018)]{jacot2018neural}
Jacot, A., Gabriel, F., and Hongler, C.
\newblock Neural tangent kernel: Convergence and generalization in neural
  networks.
\newblock In \emph{Advances in Neural Information Processing Systems}, 2018.

\bibitem[Jiang et al.(2020)]{jiang2020characterizing}
Jiang, Y., Krishnan, D., Mobahi, H., and Bengio, S.
\newblock Predicting the generalization gap in deep networks with margin
  distributions.
\newblock In \emph{International Conference on Learning Representations}, 2020.

\bibitem[Krogh and Vedelsby(1995)]{krogh1995neural}
Krogh, A. and Vedelsby, J.
\newblock Neural network ensembles, cross validation, and active learning.
\newblock In \emph{Advances in Neural Information Processing Systems}, 1995.

\bibitem[Jacot et al.(2018)]{lee2019wide}
Lee, J., Xiao, L., Schoenholz, S., Bahri, Y., Novak, R., Sohl-Dickstein, J.,
  and Pennington, J.
\newblock Wide neural networks of any depth evolve as linear models under
  gradient descent.
\newblock In \emph{Advances in Neural Information Processing Systems}, 2019.

\bibitem[Lakshminarayanan et al.(2017)]{lakshminarayanan2017simple}
Lakshminarayanan, B., Pritzel, A., and Blundell, C.
\newblock Simple and scalable predictive uncertainty estimation using deep
  ensembles.
\newblock In \emph{Advances in Neural Information Processing Systems}, 2017.

\bibitem[Ovadia et al.(2019)]{ovadia2019can}
Ovadia, Y., Fertig, E., Ren, J., Nado, Z., Sculley, D., Nowozin, S.,
  Dillon, J., Lakshminarayanan, B., and Snoek, J.
\newblock Can you trust your model's uncertainty? evaluating predictive
  uncertainty under dataset shift.
\newblock In \emph{Advances in Neural Information Processing Systems}, 2019.

\end{thebibliography}

\end{document}
